\documentclass[11pt,a4paper]{article}
\usepackage[margin=2.5cm, top=3cm, bottom=3cm]{geometry}
\usepackage{fontspec}
\usepackage{kotex}
\usepackage{booktabs}
\usepackage{array}
\usepackage{xcolor}
\usepackage{hyperref}
\usepackage{enumitem}
\usepackage{titlesec}
\usepackage{fancyhdr}
\usepackage{amsmath}
\usepackage{float}
\usepackage{caption}
\usepackage{tcolorbox}
\usepackage{tabularx}
\usepackage{colortbl}
\usepackage{setspace}

% ---------- Fonts ----------
\setmainfont[BoldFont=AppleSDGothicNeo-Bold, ItalicFont=AppleSDGothicNeo-Light]{AppleSDGothicNeo-Regular}
\setsansfont[BoldFont=AppleSDGothicNeo-Bold]{AppleSDGothicNeo-Regular}
\setmonofont{Menlo}

% ---------- Palette (house style) ----------
\definecolor{dgistnavy}{RGB}{14, 32, 80}
\definecolor{dgistblue}{RGB}{26, 60, 140}
\definecolor{dgistmid}{RGB}{60, 100, 180}
\definecolor{dgistlight}{RGB}{220, 230, 250}
\definecolor{accentgreen}{RGB}{39, 130, 67}
\definecolor{accentorange}{RGB}{210, 100, 20}
\definecolor{accentpurple}{RGB}{100, 50, 150}
\definecolor{bglight}{RGB}{248, 249, 252}
\definecolor{bgpurple}{RGB}{248, 240, 255}
\definecolor{tabhead}{RGB}{26, 60, 140}

\hypersetup{colorlinks=true, linkcolor=blue!50!black, urlcolor=blue!60!black}

% ---------- Header/Footer ----------
\pagestyle{fancy}
\fancyhf{}
\fancyhead[L]{\small\color{dgistblue}\textbf{Transformer Paper Brief}}
\fancyhead[R]{\small\color{gray} Attention Is All You Need}
\fancyfoot[C]{\small\color{gray}\thepage}
\renewcommand{\headrulewidth}{0.5pt}
\renewcommand{\headrule}{\hbox to\headwidth{\color{dgistblue}\leaders\hrule height \headrulewidth\hfill}}

% ---------- Section Styles ----------
\titleformat{\section}
  {\large\bfseries\color{dgistnavy}}
  {\colorbox{dgistblue}{\textcolor{white}{\enspace\thesection\enspace}}}
  {0.8em}{}[\vspace{2pt}\color{dgistblue}\hrule height 0.8pt\vspace{4pt}]
\titleformat{\subsection}
  {\normalsize\bfseries\color{dgistblue}}
  {\thesubsection}{0.6em}{}

\setlength{\parindent}{0pt}
\setlength{\parskip}{4pt}
\setlength{\headheight}{22pt}
\addtolength{\topmargin}{-10pt}
\onehalfspacing

% ---------- Boxes ----------
\newenvironment{keybox}[2]{%
  \begin{tcolorbox}[colback=#1, colframe=dgistblue!60,
    title={\small\bfseries\color{white}#2},
    coltitle=white, colbacktitle=dgistblue,
    left=6pt, right=6pt, top=4pt, bottom=4pt,
    boxrule=0.6pt]
}{\end{tcolorbox}}

\begin{document}

% ---------- Cover ----------
\begin{titlepage}
\pagecolor{dgistnavy}
\color{white}
\vspace*{1.4cm}
\begin{center}
{\small\color{dgistlight}\textbf{OpenClaw · Paper Design Brief}}\\[1.1cm]
{\Large\bfseries\color{white}Attention Is All You Need 간략 보고서}\\[0.5cm]
{\normalsize\color{dgistlight}A Design-Aware Summary with House LaTeX Style}\\[1.2cm]
{\Large\bfseries\color{dgistlight}2026-03-31}
\vfill
\end{center}
\end{titlepage}
\pagecolor{white}\color{black}

\tableofcontents
\newpage

\section{연구 개요}
\subsection{한 문장 요약}
\begin{keybox}{bglight}{핵심 요약}
\textit{Attention Is All You Need} (Vaswani et al., 2017)은 RNN/CNN 없이 Self-Attention만으로 시퀀스 변환을 수행하는 Transformer를 제안했고,
현대 LLM(BERT/GPT 계열)의 아키텍처 표준을 만든 전환점 논문이다.
\end{keybox}

\subsection{문제의식}
기존 RNN 기반 Seq2Seq는 (1) 장거리 의존성 학습 어려움, (2) 순차 처리로 인한 병렬화 한계, (3) 학습 효율 저하 문제가 있었다.
논문은 이 병목을 attention 중심 구조로 재정의했다.

\section{핵심 방법}
\subsection{Transformer 구조}
인코더-디코더 블록을 반복하며, 각 블록은 Multi-Head Self-Attention + Position-wise Feed-Forward + Residual/LayerNorm으로 구성된다.

\subsection{Scaled Dot-Product Attention}
\[
\mathrm{Attention}(Q,K,V)=\mathrm{softmax}\left(\frac{QK^T}{\sqrt{d_k}}\right)V
\]
$\sqrt{d_k}$ 스케일링으로 고차원에서 softmax의 불안정을 줄인다.

\subsection{Multi-Head + Positional Encoding}
여러 attention head가 서로 다른 관계를 병렬 학습하고,
positional encoding이 순서 정보를 명시적으로 주입해 RNN 없이 시퀀스 문맥을 유지한다.

\section{결과와 영향}
\begin{table}[H]
\centering\small
\begin{tabularx}{\linewidth}{|l|X|}
\hline
\rowcolor{tabhead}\textcolor{white}{\textbf{항목}} & \textcolor{white}{\textbf{핵심 내용}} \\
\hline
성능 & WMT 2014 번역 과제에서 강력한 결과를 달성하며 SOTA 수준을 제시 \\
\hline
효율 & 순차 의존을 줄여 학습 병렬화 효율을 크게 개선 \\
\hline
파급력 & BERT/GPT 등 대규모 사전학습 언어모델 아키텍처의 직접적 기반 \\
\hline
\end{tabularx}
\caption{논문의 실질적 기여}
\end{table}

\begin{keybox}{bgpurple}{실무 관점 Takeaway}
(1) Transformer 이해는 LLM 이해의 출발점이다.\\
(2) 모델 품질뿐 아니라 계산 효율(학습/추론 비용)을 함께 봐야 한다.\\
(3) Attention/Positional 설계는 실제 성능과 확장성에 직접적 영향을 준다.
\end{keybox}

\clearpage
\section{한계와 후속 연구}
대표 한계는 self-attention의 $O(n^2)$ 비용이며, 긴 문맥에서 연산/메모리 부담이 커진다.
이에 따라 Sparse/Linear Attention, FlashAttention, Long-context 최적화 연구가 이어졌다.

\section{참고 문헌}
Vaswani, A., Shazeer, N., Parmar, N., Uszkoreit, J., Jones, L., Gomez, A. N., Kaiser, \L{}., \& Polosukhin, I. (2017).
\textit{Attention Is All You Need}. NeurIPS.

\vfill
\begin{center}
\small\color{gray}\rule{0.5\linewidth}{0.4pt}\\[4pt]
OpenClaw Brief · Styled Template Edition · 2026-03-31
\end{center}

\end{document}
